\documentclass[11pt]{article}
\usepackage[margin=1in]{geometry}
\usepackage[T1]{fontenc}
\usepackage[utf8]{inputenc}
\usepackage{parskip}
\setlength{\emergencystretch}{2em}
\pagenumbering{gobble}

\begin{document}

\begin{flushright}
Zhi Huang\\
Department of Environmental Science and Engineering\\
College of Architecture and Environment, Sichuan University\\
Chengdu 610065, P.R. China\\
Institute of New Energy and Low Carbon Technology, Sichuan University\\
Chengdu 610065, P.R. China\\
Email: 525514353@qq.com
\end{flushright}

\vspace{1em}
\noindent To the Editor, \textit{Scientific Reports}

\vspace{1em}
\noindent Dear Editor,

Please find our revised manuscript, ``BioInteract: Interpretable Drug--Target
Interaction Prediction via Residue-Level Cross-Attention with Biological Prior
Knowledge'', for reconsideration in \textit{Scientific Reports}. We thank the
Editor and Reviewers for their detailed assessments. The revision is accompanied
by a clean manuscript, a marked manuscript, Supporting Information, and a
point-by-point response.

BioInteract addresses drug--target interaction prioritisation with a
residue-level cross-attention model while explicitly separating model-native
attribution from structural evidence. The revised manuscript is suited to
\textit{Scientific Reports} because it provides reproducible computational
methods, released code and artifacts, diagnostic analyses of dataset scope, and
a claim-calibrated evaluation of an AI-enabled drug-discovery workflow.

The substantive revisions include an audit of the original identifier-based
splits, exact-sequence-grouped within-Davis evaluations, and a pre-specified
strict BindingDB external evaluation. We restored the Davis baseline comparison
and component-removal results, regenerated the associated figures from the
reported values, and reconciled Figure~1 Panel~C with Table~1 by including all
seven comparison methods. We also removed the unsupported nominal ABL1
variant-comparison and hotspot-count analyses, tightened the interpretation of
attention and Grad-CAM outputs, and expanded the documentation of code,
configuration, and evaluation provenance.

For transparency, several requests that would require a new independently
benchmarked study are not represented as completed experiments. Specifically,
we report a full-length identity audit and stratified cold-target results but do
not claim a new 60\%-clustered split; we define functional-group attribution but
do not claim a matched-analogue stability test; we cite recent methods without
claiming unmatched reimplementations; we state the implementation environment
without presenting unmeasured BindingDB-scale profiling; and we do not claim a
PDBBind/LP-PDBBind overlap or structural-contact validation. In each case, the
manuscript and response letter state the scope boundary and avoid conclusions
that would require the absent evidence.

Thank you for considering our revision.

\vspace{1em}
\noindent Sincerely,\\
Zhi Huang, on behalf of all authors

\end{document}
